\vkthmsync{6}
\paragraph{\trExactTriangle{}。}\label{sub:exact-triangles}
任一抽象\trMappingCone{} $\scr{C}=\scr{C}\mathit{one}(c)$ 都自带\trCanonical{}的%
\trModule{}\trHomomorphism{} $\iota\in\hom^0_{\scr{Q}}(\scr{Y}_1,\scr{C})$ 与
$\pi\in\hom^1_{\scr{Q}}(\scr{C},\scr{Y}_0)$：按\eqref{eq:cone-definition}，它们由
$\iota^1(b_1)=(0,(-1)^{|b_1|}b_1)$、$\pi^1(b_0,b_1)=(-1)^{|b_0|-1}b_0$ 给出，
更高阶项全为零。它们在 $H(\scr{Q})$ 中构成如下三角形图：
\begin{equation}\label{eq:cone-triangle-diagram}
\begin{tikzcd}
\scr{Y}_0 \ar[rr,"{[\scr{J}^1(c)]}"] && \scr{Y}_1 \ar[dl,"{[\iota]}"] \\
& \scr{C} \ar[ul,"{[\pi]}","{[1]}"{swap}] &
\end{tikzcd}
\end{equation}
其中箭头上的记号 $[1]$ 提醒我们：相应的\trMorphism{}\trDegree{}为一。以备后文
引用，我们在此提及这些映射的另一种代表：

\begin{lemma}\label{thm:cone-alternative-representatives}
设 $e_{Y_0}$、$e_{Y_1}$ 是表示 $H(\scr{A})$ 中恒等\trMorphism{}的链。那么由
\begin{equation*}
\begin{aligned}
\tilde{\iota}^d(b,a_{d-1},\ldots,a_1)&=(0,\mu^{d+1}_{\scr{A}}(e_{Y_1},b,a_{d-1},\ldots,a_1)),\\
\tilde{\pi}^d((b_0,b_1),a_{d-1},\ldots,a_1)&=-\mu^{d+1}_{\scr{A}}(e_{Y_0},b_0,a_{d-1},\ldots,a_1)
\end{aligned}
\end{equation*}
给出的\trModule{}\trHomomorphism{} $\tilde{\iota}\in\hom^0_{\scr{Q}}(\scr{Y}_1,\scr{C})$、%
$\tilde{\pi}\in\hom^1_{\scr{Q}}(\scr{C},\scr{Y}_0)$ 分别与 $\iota$、$\pi$
\trCohomology。
\end{lemma}

\begin{proof}
我们知道\trYonedaEmbedding{}是\trCunital{}的。于是存在
$h\in\hom^{-1}_{\scr{Q}}(\scr{Y}_1,\scr{Y}_1)$，使得
$\mu^1_{\scr{Q}}(h)=e_{Y_1}-\scr{J}^1(e_{Y_1})$ 为 $\scr{Q}$ 中的严格恒等与来自
$\scr{A}$ 中恒等的某一代表元的链之差。把它取作第二个\trDirectSummand，并令第一个
为零，便得到 $k=(0,h)\in\hom^{-1}_{\scr{Q}}(\scr{Y}_1,\scr{C})$，它显然满足
$\mu^1_{\scr{Q}}(k)=\iota-\tilde{\iota}$。对 $\pi-\tilde{\pi}$ 的论证平行。
\end{proof}

我们把\textbf{\trExactTriangle}（exact triangle）定义为 $H(\scr{A})$ 中任何形如
\begin{equation}\label{eq:exact-triangle-diagram}
\begin{tikzcd}
Y_0 \ar[rr,"{[c_1]}"] && Y_1 \ar[dl,"{[c_2]}"] \\
& Y_2 \ar[ul,"{[c_3]}","{[1]}"{swap}] &
\end{tikzcd}
\end{equation}
而在\trYonedaEmbedding{}下与\eqref{eq:cone-triangle-diagram}
\trIsomorphic{}的图。这意味着：若取 $c=c_1$，则在 $H^0(\scr{Q})$ 中存在%
\trIsomorphism{} $[t]\colon Y_2\to\scr{C}=\scr{C}\mathit{one}(c)$，使得
$[\pi]\circ[t]=[\scr{J}^1(c_3)]$ 且 $[\iota]=[t]\circ[\scr{J}^1(c_2)]$。根据前文
关于 $\scr{C}\mathit{one}(c)$ 对链 $c$ 选取之依赖性的讨论，这种选取在这里其实无关
紧要（也就是说，若对某一个 $c$ 的选取能造出\trIsomorphism{} $[t]$，则对任何别的
选取亦可）。下面沿着\cref{thm:quasi-represents-characterization}的一般思路，给出
一个更具体的\trExact{}性判据：

\begin{lemma}\label{thm:exact-triangle-criterion}
三角形\eqref{eq:exact-triangle-diagram} 是\trExact{}的，当且仅当下述条件成立。
存在
\begin{equation*}
h_1\in\hom^0_{\scr{A}}(Y_1,Y_0),\quad
h_2\in\hom^0_{\scr{A}}(Y_2,Y_1),\quad
k\in\hom^{-1}_{\scr{A}}(Y_1,Y_1)
\end{equation*}
满足
\begin{equation}\label{eq:exactness-equations}
\begin{aligned}
\mu^1_{\scr{A}}(h_1)&=\mu^2_{\scr{A}}(c_3,c_2),\\
\mu^1_{\scr{A}}(h_2)&=-\mu^2_{\scr{A}}(c_1,c_3),\\
\mu^1_{\scr{A}}(k)&=-\mu^2_{\scr{A}}(c_1,h_1)+\mu^2_{\scr{A}}(h_2,c_2)
+\mu^3_{\scr{A}}(c_1,c_3,c_2)-e_{Y_1},
\end{aligned}
\end{equation}
其中 $e_{Y_1}$ 是恒等的某个链代表元；而且此外，对每个\trObject{} $X$，如下%
\trChainComplex{}必须\trAcyclic：
\begin{equation}\label{eq:exactness-complex}
\begin{gathered}
\hom_{\scr{A}}(X,Y_2)[1]\oplus\hom_{\scr{A}}(X,Y_0)[1]\oplus\hom_{\scr{A}}(X,Y_1),\\
\partial=\begin{pmatrix}
\mu^1_{\scr{A}} & 0 & 0\\
\mu^2_{\scr{A}}(c_3,\cdot) & \mu^1_{\scr{A}} & 0\\
\mu^2_{\scr{A}}(h_2,\cdot)+\mu^3_{\scr{A}}(c_1,c_3,\cdot) &
\mu^2_{\scr{A}}(c_1,\cdot) & \mu^1_{\scr{A}}
\end{pmatrix}.
\end{gathered}
\end{equation}
\end{lemma}

\begin{proof}
设上列数据都已给定。令 $\scr{C}=\scr{C}\mathit{one}(c_1)$，则 $b=(c_3,h_2)$ 是
$\scr{C}(Y_2)$ 中的\trCocycle，从而给出\trModule{}\trHomomorphism{}%
$t=\lambda_{\scr{C}}(b)\in\hom^0_{\scr{Q}}(\scr{Y}_2,\scr{C})$。\trComplex{}
\eqref{eq:exactness-complex} 是与之相应的\trChainMap{}
\eqref{eq:quasi-representing-map} 的\trMappingCone，故其\trAcyclic{}性保证 $[t]$
是 $H^0(\scr{Q})$ 中的\trIsomorphism。由\eqref{eq:lambda-multiplication-property}
及 $\scr{C}$ 的定义可见，
$[\mu^2_{\scr{Q}}(t,\scr{J}^1(c_2))]=[\lambda_{\scr{C}}(\mu^2_{\scr{C}}(b,c_2))]$
是 $(\mu^2_{\scr{A}}(c_3,c_2),\mu^2_{\scr{A}}(h_2,c_2)
+\mu^3_{\scr{A}}(c_1,c_3,c_2))\in\scr{C}(Y_1)$ 在 $\lambda_{\scr{C}}$ 下的像。按
假设\eqref{eq:exactness-equations}，它与 $(0,e_{Y_1})$ \trCohomology，而后者在
$\lambda_{\scr{C}}$ 下的像正是\cref{thm:cone-alternative-representatives} 中的%
\trModule{}\trHomomorphism{} $\tilde{\iota}$。于是
\begin{equation*}
[t]=[\mu^2_{\scr{Q}}(t,\scr{J}^1(c_2))].
\end{equation*}
其次，利用\eqref{eq:lambda-second-property} 可见，在
$\Hom^1_{H(\scr{Q})}(\scr{Y}_2,\scr{Y}_0)$ 中
\begin{equation*}
[\mu^2_{\scr{Q}}(\pi,t)]=[\lambda_{\scr{Y}_0}(\pi^1(b))]
=[\lambda_{\scr{Y}_0}(c_3)]=[\scr{J}^1(c_3)].
\end{equation*}
至此，\trExactTriangle{}的各项要求已全部列出。

反方向基本上就是把上述论证倒着读一遍。任何\trMorphism{}%
$[t]\colon\scr{Y}_2\to\scr{C}$ 都经由 $\lambda_{\scr{C}}$ 来自某个元素
$b\in\scr{C}(Y_2)$。由于
$[\scr{J}^1(\pi^1(b))]=[\mu^2_{\scr{Q}}(\pi,t)]=[\scr{J}^1(c_3)]$，$b$ 的第一个
分量与 $c_3$ \trCohomology。加上一个适当的\trCoboundary{}之后，不妨设
$b=(c_3,h_2)$。接下来，
$[\lambda_{\scr{C}}(0,e_{Y_1})]=[\tilde{\iota}]
=[\mu^2_{\scr{Q}}(t,\scr{J}^1(c_2))]
=[\lambda_{\scr{C}}(\mu^2_{\scr{C}}(b,c_2))]
=[(\mu^2_{\scr{A}}(c_3,c_2),\mu^2_{\scr{A}}(h_2,c_2)
+\mu^3_{\scr{A}}(c_1,c_3,c_2))]$。在这两个\trCycle{}之间选取一个\trBoundingCochain，便给出
其余数据 $(k,h_1)$。
\end{proof}
